%-------------------------------------------------
% FileName: chap-7.tex
% Description: 第7章 总结与展望
%-------------------------------------------------

\chapter{总结与展望}

\section{研究总结}

本文用 Java Web 技术实现了一个体育馆管理系统。本系统以校园体育馆里较为常见的场地预约、器材租赁为依托，分别对普通用户、管理员两大类用户的使用需求进行分析，最后完成了前后端分离的架构设计、数据库表设计、主要功能实现及测试工作。本文在整体开发的过程中并没有将体育馆管理抽象为一个巨大的平台，而是选取了预约、订单、器材、后台维护等几个具体环节来开展工作，保证系统可以形成可以演示和测试的闭环。

用户端的功能有注册登录、场地列表、预约下单、我的订单、器材租借和个人记录查询等。预约流程中，系统会根据场地开放时间及当天已占用时段生成可选时间段，提交订单前后端都会进行校验；订单生成之后可以进行支付或者取消，状态改变会反映到个人订单列表上。管理端实现了场地管理、订单查看、器材维护、用户角色管理以及数据看板等主要功能，管理员能够对基础数据进行维护，也可以在后台首页上看到目前系统的大概运行状况。

在技术实现方面，前端采用 Vue3、Vite 和 Element Plus，后端采用 Spring Boot、MyBatis-Plus 和 MySQL。系统用 JWT 保存登录后的身份信息，在后端接口中按照用户编号和角色来限制访问范围；统一返回对象、全局异常处理使前端可以以类似的方式去处理接口的结果。测试阶段主要是针对登录认证、场地预约、订单状态流转、库存扣减、归还处理和管理员权限等流程进行测试，结果证明本系统可以完成毕业设计所要达到的功能要求。

但是系统存在一定的不足。由于课题时长受限以及开发范围的制约，目前支付功能仍然使用状态模拟的方式，并未与真实的支付通道对接；密码存储、日志审计、接口限流等安全保障措施不齐全，预约冲突依靠业务查询和事务处理解决，在更高的并发下仍需加强数据库约束或者使用锁。这些不足可以明确系统边界的把握情况，并给后面改进建设给予线索。

\section{后续展望}

之后可以从业务完整性上进行完善。真实使用时预约支付不应该只是订单状态的切换，还可以接入支付接口或者校园卡系统；预约成功、订单取消、器材归还超时等节点也可以增加短信、邮件或者站内通知提醒。普通用户一般是通过手机进行预订，因此前段可以更加深入地改善小屏幕上表格所占用的部分，在减少用户阅读和使用时间的基础上进一步提高效率。

安全性也是后面需要加强的部分。目前系统适合课程演示，但是要接近真实部署的话，密码应该用加盐哈希存储，登录失败次数也应该有限制，关键管理操作也应当记录日志，管理员权限也可以从简单的 admin/user 二分扩展为更细的权限项。将场地维护、订单查看、用户角色管理拆分，避免使用同一个管理员账号具有太大权限。

就性能以及数据分析而言，系统能渐渐地把统计逻辑从前端迁移至后端聚合接口之中。随着订单数量的增加，后台首页不能再一次性读取大量的订单然后交给前端去计算，而应该由后端按日期、场地类型或者状态直接返回统计结果。还可以添加场地使用率、热门时间段、器材借还情况等分析项，利于管理员对场地进行调度安排。

最后，学校如果有统一的身份认证、校园卡或者门禁系统等，那么本系统就可以和这些平台进行对接。因此用户不必另设账号，管理员亦可将预约、支付、入馆核验等关联起来。由于本科毕业设计开发条件所限，本文并没有实现这些扩展功能，但是从系统的模块划分来看，在认证、支付、统计模块上继续进行迭代是可行的。
